\documentclass[./测控三级项目.tex]{subfiles} % 使用 subfiles 文档类，指定主文档

\begin{document}
	\subsection{题目重述}
		\begin{compactenum}
			\item 设计一个同相一阶有源低通滤波器，分别构建仿真电路和实测电路，滤除题6[1]d)中的高频信号。
			\item 设计一个同相一阶有源高通滤波器，分别构建仿真电路和实测电路，滤除题6[1]d)中的低频信号。
			\item 描述同相一阶有源滤波器与一阶无源RC滤波器的性能差异。
			\item 搭建实测电路，分别为上述低通实测电路引入一个20mV正弦信号，调整信号频率，在0.1Hz-1MHz 范围内测试30个以上关键频率点的输出波形，记录输出信号幅值变化，用MATLAB描点绘制Bode图，用术语描述滤波器的频率特性源。
		\end{compactenum}
	\subsection{题目分析与实施方案}
		\subsubsection{第一问}
			第一问要求设计同相有源低通滤波器，最常见的有源一阶低通滤波器是，因此考虑在本题使用这个电路。\par
			电路确定后，通过计算无源反馈网络的传递函数，得到对应相频特性函数。将需要滤除的频率代入特性曲线的截止频率，便可以推算出对应电路参数，并根据实际情况，选用相关精度的标准阻值电阻。\par
			然后，使用Multisim软件仿真对应电路，将题6[1]d)中的输出信号输入搭建的德玻滤波器上，并连接示波器。观察示波器波形，判断高频信号是否被滤除。
		\subsubsection{第二问}
			第二问要求设计同相有源高通滤波器，最常见的有源高通滤波器是一阶微分器电路，因此考虑在本题使用这个电路。\par
			电路确定后，通过计算无源反馈网络的传递函数，得到对应相频特性函数。将需要滤除的频率代入特性曲线的截止频率，便可以推算出对应电路参数，并根据实际情况，选用相关精度的标准阻值电阻。\par
			然后，使用Multisim软件仿真对应电路，将题6[1]d)中的输出信号输入搭建的一阶微分器电路上，并连接示波器。观察示波器波形，判断低频信号是否被滤除。
		\subsubsection{第三问}
			第三问要求描述同相一阶有源滤波器与一阶无源RC滤波器的性能差异，可根据本题设计的同相一阶有源滤波器产生的结果与上一题种设计一阶无源RC滤波器进行比较，最终得出可靠的结论。
		\subsubsection{第四问}
			根据第一问搭建的电路，向有源低通滤波器等对数频率，分别测试输出信号的幅值变化三十组。\par
			使用MATLAB的plot函数，描点绘制出测试得到的参数，以得到相频伯德图曲线。\par
			根据绘制的伯德图曲线，分析该RC无源低通滤波器的频率特性。
	
	\subsection{理论分析与仿真}
		\subsubsection{第一问}
			图\ref{fig:同相一阶有源低通滤波器电路原理图}是同相一阶有源低通滤波器电路原理图。
			\begin{figure}[h]
				\centering
				\includegraphics[width=0.45\textwidth]{同相一阶有源低通滤波器电路原理图}
				\caption{同相一阶有源低通滤波器电路原理图}
				\label{fig:同相一阶有源低通滤波器电路原理图}
			\end{figure}\par
			根据电路计算可得：
			\begin{gather}
				U_o=(1+\frac{R_F}{R_1})\frac{1}{RCs+1}U_i
			\end{gather}\par
			由此可得：
			\begin{gather}
				G(s)=\frac{u_o}{u_i}=(1+\frac{R_F}{R_1})\frac{1}{RCs+1}
				\label{eq:同相一阶有源低通滤波器传递函数}
			\end{gather}\par
			令$s=j\omega$，并将其表示为归一化形式，可得：
			\begin{align}
				G(j\omega)&=(1+\frac{R_F}{R_1})\frac{1}{j\omega RC+1}\\
				&=H_0\frac{1}{1+j\omega/\omega_0}
			\end{align}\par
			因此归一化频率$\omega_0$和归一化增益$H_0$，满足：
			\begin{gather}
				\begin{cases}
					\omega_0 =\frac{1}{RC}=2\pi f_0
					\\
					H_0=1+\frac{R_F}{R_1}
				\end{cases}
				\label{eq:同相一阶有源低通滤波器截止频率}
			\end{gather}\par
			以上理论完备，在Multisim软件中进行仿真实验。\par
			由于需要滤除的高频信号为5kHz，需要保留的低频信号为2kHz。因此，设置截止频率为$f_0=4$kHz。设置电容容值为$C=1$nF(Standeard Value,$5\%$)\par
			根据式\ref{eq:同相一阶有源低通滤波器截止频率}，电阻阻值$R=\frac{1}{2\pi Cf_0}=39.788$k$\Omega$，与标准电阻值表比对后，实际选用电阻阻值为$R=40.2$k$\Omega$(Standeard Value,$1\%$)的电阻。\par
			同时，我们希望使用同相放大器对低频信号进行补偿，即当$f_L=2$kHz时，增益为1：
			\begin{gather}
				A(\omega)=\lvert G(j\omega) \rvert =(1+\frac{R_F}{R_1})\frac{1}{\sqrt{1+(RC\omega_L)^2}}
			\end{gather}\par
			将$\omega_L=2\pi f_L=4000\pi$、$R=40.2$k$\Omega$、$C=1$nF代入计算，整理后得到：
			\begin{gather}
				R_F=0.1203547688R_1
			\end{gather}\par
			设置电阻$R_1$的阻值为$R_1=150$k$\Omega$(Standeard Value,$1\%$)，则$R_F$阻值应为$R_F=18.05321532$k$\Omega$，与标准电阻值表比对后，实际选用电阻阻值为$R_F=18.2$k$\Omega$(Standeard Value,$1\%$)的电阻。\par
			将要求的信号接入该电路，得到仿真电路图，如图\ref{fig:同相一阶有源低通滤波器仿真原理图}所示。
			\begin{figure}[H]
				\centering
				\includegraphics[width=0.45\textwidth]{./8.第八题示波器现象/一阶有源低通滤波器仿真原理图}
				\caption{同相一阶有源低通滤波器仿真原理图}
				\label{fig:同相一阶有源低通滤波器仿真原理图}
			\end{figure}\par
			对输入和输出信号做傅里叶变换，得到对应频谱曲线，如图\ref{fig:仿真同相一阶有源低通滤波器频谱曲线}所示。
			\begin{figure}[h]
				\centering
				\includegraphics[width=0.45\textwidth]{./8.第八题示波器现象/仿真一阶有源低通滤波器频谱曲线}
				\caption{仿真同相一阶有源低通滤波器频谱曲线}
				\label{fig:仿真同相一阶有源低通滤波器频谱曲线}
			\end{figure}\par
			对于低频信号（2kHz正弦信号），幅值从0.399455降到了0.399694，可求得增益为$A=-20\lg(\frac{0.399455}{0.399694})=0.005195345$dB；对于高频信号（5kHz正弦信号），幅值从0.0594916降到了0.0411997，可求得增益为$A=-20\lg(\frac{0.0594916}{0.0411997})=-3.191231911$dB。\par
			这说明图\ref{fig:同相一阶有源低通滤波器电路原理图}的低通滤波器可以对高频信号进行滤波操作。
		\subsubsection{第二问}
			图\ref{fig:同相一阶有源高通滤波器电路原理图}是同相一阶有源高通滤波器电路原理图。
			\begin{figure}[h]
				\centering
				\includegraphics[width=0.45\textwidth]{同相一阶有源高通滤波器电路原理图}
				\caption{同相一阶有源高通滤波器电路原理图}
				\label{fig:同相一阶有源高通滤波器电路原理图}
			\end{figure}\par
			根据电路计算可得：
			\begin{gather}
				U_o=(1+\frac{R_F}{R_1})\frac{R}{\frac{1}{Cs}+R}U_i=(1+\frac{R_F}{R_1})\frac{RCs}{RCs+1}U_i
			\end{gather}\par
			由此可得：
			\begin{gather}
				G(s)=\frac{u_o}{u_i}=(1+\frac{R_F}{R_1})\frac{RCs}{RCs+1}
				\label{eq:同相一阶有源高通滤波器传递函数}
			\end{gather}\par
			令$s=j\omega$，并将其表示为归一化形式，可得：
			\begin{align}
				G(j\omega)&=(1+\frac{R_F}{R_1})\frac{j\omega RC}{j\omega RC+1}\\
				&=H_0\frac{j\omega/\omega_0}{1+j\omega/\omega_0}
			\end{align}\par
			因此归一化频率$\omega_0$和归一化增益$H_0$，满足：
			\begin{gather}
				\begin{cases}
					\omega_0 =\frac{1}{RC}=2\pi f_0
					\\
					H_0=1+\frac{R_F}{R_1}
				\end{cases}
				\label{eq:同相一阶有源高通滤波器截止频率}
			\end{gather}\par
			以上理论完备，在Multisim软件中进行仿真实验。\par
			由于需要滤除的低频信号为2kHz，需要保留的高频信号为5kHz。因此，设置截止频率为$f_0=4$kHz。设置电容容值为$C=1$nF(Standeard Value,$5\%$)\par
			根据式\ref{eq:同相一阶有源高通滤波器截止频率}，电阻阻值$R=\frac{1}{2\pi Cf}=39.788$k$\Omega$，与标准电阻值表比对后，实际选用电阻阻值为$R=40.2$k$\Omega$(Standeard Value,$1\%$)的电阻。\par
			同时，我们希望使用同相放大器对高频信号进行补偿，即当$f_H=5$kHz时，增益为1：
			\begin{gather}
				A(\omega)=\lvert G(j\omega) \rvert =(1+\frac{R_F}{R_1})\frac{RC\omega_H}{\sqrt{1+(RC\omega_H)^2}}
			\end{gather}\par
			将$\omega_H=2\pi f_L=10000\pi$、$R=40.2$k$\Omega$、$C=1$nF代入计算，整理后得到：
			\begin{gather}
				R_F=0.2755281273R_1
			\end{gather}\par
			设置电阻$R_1$的阻值为$R_1=110$k$\Omega$(Standeard Value,$1\%$)，则$R_F$阻值应为$R_F=30.308094$k$\Omega$，与标准电阻值表比对后，实际选用电阻阻值为$R_F=30.1$k$\Omega$(Standeard Value,$1\%$)的电阻。\par
			将要求的信号接入该电路，得到仿真电路图，如图\ref{fig:同相一阶有源高通滤波器仿真原理图}所示。
			\begin{figure}[H]
				\centering
				\includegraphics[width=0.45\textwidth]{./8.第八题示波器现象/一阶有源高通滤波器仿真原理图}
				\caption{同相一阶有源高通滤波器仿真原理图}
				\label{fig:同相一阶有源高通滤波器仿真原理图}
			\end{figure}\par
			对输入和输出信号做傅里叶变换，得到对应频谱曲线，如图\ref{fig:仿真同相一阶有源高通滤波器频谱曲线}所示。
			\begin{figure}[H]
				\centering
				\includegraphics[width=0.45\textwidth]{./8.第八题示波器现象/仿真一阶有源高通滤波器频谱曲线}
				\caption{仿真同相一阶有源高通滤波器频谱曲线}
				\label{fig:仿真同相一阶有源高通滤波器频谱曲线}
			\end{figure}\par
			对于低频信号（2kHz正弦信号），幅值从0.399461降到了0.229812，可求得增益为$A=-20\lg(\frac{0.399461}{0.229812})=-4.802033651$dB；对于高频信号（5kHz正弦信号），幅值从0.0594922降到了0.0597016，可求得增益为$A=-20\lg(\frac{0.0594922}{0.0597016})=0.0305188$dB。
			这说明图\ref{fig:同相一阶有源高通滤波器电路原理图}的高通滤波器可以对低频信号进行滤波操作。
		\subsubsection{第三问}
			一阶无源RC滤波器和一阶有源滤波器在性能上的差异主要体现在以下几个方面：\par
			\textbf{增益和放大能力}\qquad
			一阶无源RC滤波器没有放大功能，其增益通常小于1，因为它是由无源元件（电阻和电容）组成，不能放大信号；一阶有源滤波器由于包含运算放大器等有源元件，可以实现信号的放大，其增益可以通过调整电路参数（如反馈电阻）来控制，通常大于1。\par
			\textbf{负载效应}\qquad
			一阶无源RC滤波器其滤波特性受负载影响较大，因为无源元件本身没有隔离和缓冲作用，负载阻抗的变化会影响滤波器的截止频率和通带增益；一阶有源滤波器运算放大器的高输入阻抗和低输出阻抗特性使得滤波器的输出端不受负载影响，能够保持稳定的滤波特性。\par
			\textbf{频率响应}\qquad
			一阶无源RC滤波器其频率响应相对平缓，截止频率附近的衰减较慢，通带内信号有一定的能量损耗；一阶有源滤波器可以实现更陡峭的滤波特性，即在截止频率附近信号衰减更快，能够更有效地滤除不需要的频率成分。\par
			\textbf{功耗}\qquad
			一阶无源RC滤波器由于没有有源元件，因此功耗低；一阶有源滤波器由于包含有源元件，需要外部电源供电，因此功耗较高。
			
	\subsection{实际电路测试结果与数据分析}
		\subsubsection{第一问}
			根据上述电路分析，选用上述计算对应的电阻和电容，按照图\ref{fig:同相一阶有源低通滤波器电路原理图}的电路焊接在PCB板上\par
			焊接后，经检查分立元件位置端正、放置平整、焊点表面光滑、无凸起毛刺、无虚焊，周围无助焊剂残留；多引脚器件无短路，焊点均匀一致。符合测试要求。\par
			将信号发生器的三路信号分别加载在三路加法器的信号输入端，将示波器通道CH1接入同相一阶有源低通滤波器的输出端。添加FFT算子计算输出通道CH1的幅频曲线。\par
			将电路连接完毕后，先打开线性电源输出开关，随后根据《测控仪器电路项目作业指导书》中要求，加载对应正弦信号。\par
			同相一阶有源低通滤波器实验实测对应的频域曲线，如图\ref{fig:同相一阶有源低通滤波器滤波前后信号幅频曲线}所示。
			\begin{figure}[H]
				\centering
				\begin{subfigure}{0.4\linewidth}
					\centering
					\includegraphics[width=0.75\linewidth]{./8.第八题示波器现象/UpoPrtsc-001}
					\caption{同相一阶有源低通滤波器输入信号幅频曲线}
					\label{fig:同相一阶有源低通滤波器输入信号幅频曲线}
				\end{subfigure}
				\begin{subfigure}{0.4\linewidth}
					\centering
					\includegraphics[width=0.75\linewidth]{./8.第八题示波器现象/UpoPrtsc-002}
					\caption{同相一阶有源低通滤波器输出信号幅频曲线}
					\label{fig:同相一阶有源低通滤波器输出信号幅频曲线}
				\end{subfigure}
				\caption{同相一阶有源低通滤波器滤波前后信号幅频曲线}
				\label{fig:同相一阶有源低通滤波器滤波前后信号幅频曲线}
			\end{figure}\par
			对比图\ref{fig:同相一阶有源低通滤波器输入信号幅频曲线}和图\ref{fig:同相一阶有源低通滤波器输出信号幅频曲线}，可以发现低频的2kHz幅值基本不变；高频的5kHz信号幅值有所减小。这说明了同相一阶有源低通滤波器对于高频交流信号的幅值会有所抑制，且可以通过反相端的纯电阻反馈网络进行低频信号幅值的补偿。
		\subsubsection{第二问}
			根据上述电路分析，选用上述计算对应的电阻和电容，按照图\ref{fig:同相一阶有源高通滤波器电路原理图}的电路焊接在PCB板上\par
			焊接后，经检查分立元件位置端正、放置平整、焊点表面光滑、无凸起毛刺、无虚焊，周围无助焊剂残留；多引脚器件无短路，焊点均匀一致。符合测试要求。\par
			将信号发生器的三路信号分别加载在三路加法器的信号输入端，将示波器通道CH1接入同相一阶有源高通滤波器的输出端。添加FFT算子计算输出通道CH1的幅频曲线。\par
			将电路连接完毕后，先打开线性电源输出开关，随后根据《测控仪器电路项目作业指导书》中要求，加载对应正弦信号。\par
			同相一阶有源高通滤波器实验实测对应的频域曲线，如图\ref{fig:同相一阶有源高通滤波器滤波前后信号幅频曲线}所示。
			\begin{figure}[H]
				\centering
				\begin{subfigure}{0.4\linewidth}
					\centering
					\includegraphics[width=0.75\linewidth]{./8.第八题示波器现象/UpoPrtsc-001}
					\caption{同相一阶有源高通滤波器输入信号幅频曲线}
					\label{fig:同相一阶有源高通滤波器输入信号幅频曲线}
				\end{subfigure}
				\begin{subfigure}{0.4\linewidth}
					\centering
					\includegraphics[width=0.75\linewidth]{./8.第八题示波器现象/UpoPrtsc-003}
					\caption{同相一阶有源高通滤波器输出信号幅频曲线}
					\label{fig:同相一阶有源高通滤波器输出信号幅频曲线}
				\end{subfigure}
				\caption{同相一阶有源高通滤波器滤波前后信号幅频曲线}
				\label{fig:同相一阶有源高通滤波器滤波前后信号幅频曲线}
			\end{figure}\par
			对比图\ref{fig:同相一阶有源高通滤波器输入信号幅频曲线}和图\ref{fig:同相一阶有源高通滤波器输出信号幅频曲线}，可以发现高频的5kHz幅值基本不变；低频的2kHz信号幅值有所减小。这说明了同相一阶有源高通滤波器对于低频交流信号的幅值会有所抑制，且可以通过反相端的纯电阻反馈网络进行信号幅值的补偿。
		\subsubsection{第四问}
			按题目要求引入幅值为20mV的正弦信号，并调整信号频率，在0.1Hz-1MHz区间测试30个以上关键频率点，测试结果如表\ref{tab:不同频率下同相一阶有源低通滤波器增益表}所示。
			\begin{table}[h!]
				\centering
				\caption{不同频率下同相一阶有源低通滤波器增益表}
				\label{tab:不同频率下同相一阶有源低通滤波器增益表}
				\begin{tabular}{|c|c|c|c|c|c|c|c|}
					\hline
					频率（Hz） & 增益  & 频率（Hz） & 增益 & 频率（Hz） & 增益 & 频率（Hz） & 增益\\
					\hline
					0.1&1.13&10&1.13&1000&1.09&100000&0.05\\
					\hline
					0.3&1.13&30&1.13&3000&0.9&300000&0.02\\
					\hline
					0.5&1.13&50&1.13&5000&0.7&500000&0.01\\
					\hline
					0.7&1.13&70&1.13&7000&0.56&700000&0.01\\
					\hline
					0.9&1.13&90&1.13&9000&0.46&900000&0.01\\
					\hline
					1&1.13&100&1.13&10000&0.42&1000000&0.01\\
					\hline
					3&1.13&300&1.12&30000&0.15&&\\
					\hline
					5&1.13&500&1.12&50000&0.09&&\\
					\hline
					7&1.13&700&1.11&70000&0.07&&\\
					\hline
					9&1.13&900&1.1&90000&0.05&&\\
					\hline
				\end{tabular}
			\end{table}	\par
			将表\ref{tab:不同频率下同相一阶有源低通滤波器增益表}中的数据输入MATLAB中进行Bode图生成，其图像如下图\ref{fig:同相一阶有源低通滤波器测试伯德图}所示。
			\begin{figure}[h]
				\centering
				\includegraphics[width=0.7\textwidth]{./8.第八题示波器现象/同相一阶有源低通滤波器测试伯德图}
				\caption{测试生成同相一阶有源低通滤波器伯德图}
				\label{fig:同相一阶有源低通滤波器测试伯德图}
			\end{figure}\par
			根据实际选取电阻电容值参数，可以直接得到对应传递函数为$G(s)=\frac{1+\frac{18.2}{150}}{4.02\times  10^{-5}s+1}$。利用MATLAB软件，得到该传递函数的伯德图如图\ref{fig:同相一阶有源低通滤波器实际伯德图}所示。
			\begin{figure}[H]
				\centering
				\includegraphics[width=0.7\textwidth]{./8.第八题示波器现象/同相一阶有源低通滤波器实际伯德图}
				\caption{理论生成同相一阶有源低通滤波器伯德图}
				\label{fig:同相一阶有源低通滤波器实际伯德图}
			\end{figure}\par
			对比图\ref{fig:同相一阶有源低通滤波器测试伯德图}和图\ref{fig:同相一阶有源低通滤波器实际伯德图}，中的相频曲线，可以发现除了1kHz-30kHz除下降曲线前者不如后者光滑外，其余基本一致，这说明实际测试结果符合理论推断。\par
			同时也注意到，在1kHz-30kHz范围内，增益大幅变化，从1下降到了接近0的位置，这说明低通滤波器在该频段具有抑制高频信号的特点。
	\subsection{结论与收获}
		\subsubsection{结论}
			8.1经过实验验证，图\ref{fig:同相一阶有源低通滤波器电路原理图}所设计的电路基本满足同向一阶低通滤波器要求，可以对高频信号进行过滤。\par
			8.2经过实验验证，图\ref{fig:同相一阶有源高通滤波器电路原理图}所设计的电路基本满足同向一阶高通滤波器要求，可以对低频信号进行过滤。\par
			8.3同相一阶有源滤波器与一阶无源RC滤波器的性能差异：同相一阶有源滤波器内有增益，能对信号进行放大，而一阶无源 RC 滤波器无此功能；在负载效应方面，同相一阶有源滤波器因运算放大器的高输入阻抗和低输出阻抗特性，带负载能力强，负载对其频率特性影响小，一阶无源 RC滤波器则因输出阻抗较高，负载会改变其频率特性；同相一阶有源滤波器虽可放大信号，但是也容易引起失真，一阶无源 RC 滤波器则无失真，但是对信号有衰减作用。\par
			8.4根据绘制出的 Bode 图以及实验数据可知，该同相一阶有源低通滤波器具有典型的低通特性。在低频段，输出信号幅值与输入信号幅值相近，随着频率逐渐升高，输出信号幅值开始缓慢下降，当频率接近截止频率时，下降速率加快。在截止频率处，输出信号幅值下降到输入信号幅值的约 0.707 倍（即 -3dB 点）。超过截止频率后，输出信号幅值按照 -20dB / 十倍频的速率持续下降，这表明滤波器对高频信号具有良好的衰减作用，能够有效地滤除高于截止频率的高频成分，从而实现低通滤波功能，使低频信号顺利通过，高频信号得到抑制，进而满足对输入信号进行低通滤波处理的需求，可应用于如音频信号处理中去除高频噪声、传感器信号采集时滤除高频干扰等场景，以提高信号的质量和后续处理的准确性。
		\subsubsection{收获}
			\textbf{游寅龙}：\par
			理论分析层面：主要负责了计算后仿真电路图的搭建和对计算的结果进行了验证，经过这些学习，对于同向一阶滤波器有了更加全面的认识，能够更熟练地使用知识解决问题。\par
			动手实践层面：完成了20mV 正弦信号的测试和Matlab仿真，能够更加熟练的使用信号源和示波器，对于Matlab软件的掌握也更进一步，做到了理论实践相结合。\par
			精神层面：通过与小组协作不断纠正错误和查找问题，培养了不怕挫折的韧性，也时刻激励着自己不断实事求是。\par
			\textbf{郭嘉媛}：\par
			理论分析层面：主要负责了同相一阶有源低通滤波器的设计和计算，明晰了低通滤波器中电容与电阻构成的 RC 网络对高频信号的衰减作用，对于知识掌握更进一步。\par
			动手实践层面：交替完成了同相一阶有源高通滤波器的焊接与测试，能够因地制宜，面对不熟悉的电路图也能够依靠自身找到解决问题的方法，提高了面对未知困难的能力。\par
			精神层面：培养了严谨的科学态度和勇于探索的精神，面对困难时能够根据多种现有的资料给出合理的解决方案，能够充分利用已有资源解决问题。\par
			\textbf{郑嘉宁}：\par
			理论分析层面：主要负责了同相一阶有源高通滤波器的设计和计算，了解了如何根据电容和电阻的值来影响电路，对于高通滤波器的知识更加熟练。\par
			动手实践层面：交替完成了同相一阶有源低通滤波器的焊接与测试，能够根据电路图合理的进行布线和焊接，焊接速度极大提升。\par
			精神层面：在设计和实测时总会遇到许多困难，能够有严谨耐心的态度去排查和解决问题，面对解决不了的问题能够发挥团队精神来解决，增强了团队凝聚力。
\end{document}